\documentclass[11pt]{amsart}
\usepackage{a4wide,amsfonts,amssymb,stmaryrd,amscd,amsmath,latexsym,amsbsy}
%\documentstyle{article}
\theoremstyle{plain}
\newtheorem{theorem}{Theorem}[section]
\newtheorem{thm}{Theorem}[section]
\newtheorem*{thrm}{Theorem}
\newtheorem{prop}[theorem]{Proposition}
\newtheorem{pro}[theorem]{Proposition}
%\newtheorem{defi}{Definition}
\newtheorem{lem}[theorem]{Lemma}
\newtheorem{cor}[theorem]{Corollary}
\newtheorem{eg}{{Example}}
\newtheorem{ex}{Exercise}
\newtheorem*{sol}{Solution}
\newtheorem*{LRrule}{Littlewood-Richardson rule}
\newtheorem*{conj}{Conjecture}
%\redefine{thebibliography}{Publications}
\theoremstyle{definition}
\newtheorem*{defi}{Definition}
\theoremstyle{remark}
\newtheorem{rema}{Remark}[section]
\newcommand{\CC}{\mathbb C}
\newcommand{\ZZ}{\mathbb Z}
\newcommand{\solu}[1]{\begin{sol}{\bf (\ref{#1})}}
\newcommand{\dd}{\partial}
\newcommand{\gotg}{\mathfrak g}
\newcommand{\gotk}{\mathfrak k}
\pagestyle{plain}
\begin{document}
\section*{\bf{\huge Alexei Oblomkov}}


\section*{\Large\underline{Recent Publications}}
\noindent\hangindent=0.8 cm 1. (with V. Shende) 
{\it The Hilbert scheme of a plane curve singularity and the HOMFLY polynomial of its link}, arxive:1003.1568,
  Duke Math. J. 161 (2012), no. 7, 1271--1303.

\noindent\hangindent=0.8 cm 2. (with V. Shende, J. Rassmusen)
{\it The Hilbert scheme of a plane curve singularity and HOMFLY homology of its link},  Geom. Topol. 22 (2018), no. 2, 645--691, arXive:1201.2115.


\noindent\hangindent=0.8 cm 3. (with E. Gorsky, J. Rassmusen) {\it On stable Khovanov homology of torus knots},    arXiv:1206.2226,  Exp. Math. 22 (2013), no. 3, 265--281.

\noindent\hangindent=0.8 cm 4. (with E. Gorsky, V. Shende, J. Rassmusen), {\it Torus knots and Rational DAHA},   arXiv:1207.4523,  Duke Math. J. 163, (2014), no. 15, 325--401. 

\noindent\hangindent=0.8 cm 5. (with Z. Yun) {\it Geometric Representations of graded and rational Cherednik algebras},   arXiv:1407.5685,
 Adv. Math. 292 (2016), 601--706.

\noindent\hangindent=0.8 cm 6. (with S. Nawata) {\it Lectures on knot homology},   arXiv:1407.5685,   Proceedings of "Workshop on Physics and Mathematics of Link Homology", at Centre de Recherches Mathematiques, Universite de Montreal,  137--177, Contemp. Math., 680, Amer. Math. Soc., Providence, RI, 2016.


\noindent\hangindent=0.8 cm 7. (with L. Rozansky) {\it Knot Homology and sheaves on the Hilbert scheme of points on the plane },
Selecta Math. (N.S.) 24 (2018), no. 3, 2351--2454,  arXiv:1608.03227.



\noindent\hangindent=0.8 cm 8. (with L. Rozansky) {\it Affine Braid group, JM elements and knot homology},
Transform. Groups 24 (2019), no. 2, 531--544. arXiv:1702.03569.


\noindent\hangindent=0.8 cm 9. (with L. Rozansky) {\it HOMFLYPT homology of Coxeter links},  arXiv:1706.00124.




\noindent\hangindent=0.8 cm 10. (with  Z. Yun) {\it The cohomology ring of certain compactified Jacobians},  arXiv:1710.05391. 


\noindent\hangindent=0.8 cm 11. (with A. Okounkov, R. Pandharipande) {\it  GW/PT descendent correspondence via vertex operators},
accepted in Communications in Mathematical Physics, arXiv:1806.00714.



\noindent\hangindent=0.8 cm 12. (with L. Rozansky)
{\it A categorification of a cyclotomic Hecke algebra}, arXiv:1801.06201.

\noindent\hangindent=0.8 cm 13. (with E. Carlsson)
{\it Affine Schubert calculus and double coinvariants}, arXiv:1801.09033.

\noindent\hangindent=0.8 cm 14. (with L. Kamenova, G. Mongardi) {\it Symplectic involutions of \(K3^{[n]}\) type and Kummer \(n\)
  type manifolds}, arXiv:1809.02810.

\noindent\hangindent=0.8 cm 15. (with L. Rozansky), {\it Categorical Chern character and braid groups}, arXiv:1811.03257.

\noindent\hangindent=0.8 cm 16. (with L. Rozansky), {\it 3D TQFT and HOMFLYPT homology}, arXiv:1812.06340.

\noindent\hangindent=0.8 cm 17. {\it EGL formula for DT/PT theory of local curves}, arXiv:1901.03014.

\noindent\hangindent=0.8 cm 18. {\it Notes on matrix factorizations and knot homology}, arXiv:1901.04052.

\noindent\hangindent=0.8 cm 19. (with L. Rozansky), {\it Dualizable link homology }, arXiv:1905.06511.

\noindent\hangindent=0.8 cm 20. (with N. Bottman), {\it 
A compactification of the moduli space of marked vertical lines in $\mathbb{C}^2$, } arXiv:1910.02037.



\end{document}


%%% Local Variables:
%%% mode: latex
%%% TeX-master: t
%%% End:
